\documentclass[12pt, a4paper]{article}
\usepackage[utf8]{inputenc}
\usepackage[T1]{fontenc}
\usepackage{amsmath, amssymb}
\usepackage{graphicx}
\usepackage{float}
\usepackage{natbib}   % لنظام Harvard (author-year)
\usepackage{geometry}
\geometry{margin=1in}
\usepackage{hyperref}
\hypersetup{
    colorlinks=true,
    citecolor=blue,
    urlcolor=blue,
}

\title{A Hybrid Quantum-Classical Optimization Framework for Well Placement in Oil Reservoirs Using QUBO Modeling and QAOA}
\author{
    Tareq Majeed Al-Karimi\thanks{Ministry of Oil, Iraq} \\
    \and
    Rasha B. Yousif\thanks{Department of Computer Science, College of Computer Science and Mathematics, University of Thi-Qar, Iraq}
}
\date{}

\begin{document}

\maketitle

\begin{abstract}
Optimal well placement remains one of the most important decision problems in field development planning due to its direct impact on production performance, reservoir drainage efficiency, and project economics. Traditional optimization approaches frequently encounter scalability limitations when the number of candidate well locations increases, particularly when the objective combines production gain and inter-well interaction penalties.

This study presents a hybrid quantum-classical framework for well placement optimization using Quadratic Unconstrained Binary Optimization (QUBO) modeling and the Quantum Approximate Optimization Algorithm (QAOA). The proposed framework reformulates the reservoir decision problem into a binary optimization structure suitable for quantum-inspired computation. The methodology is benchmarked against classical methods including exact optimization, simulated annealing, and random search using synthetic scenarios, and is further extended to a realistic field-inspired case derived from the publicly available Volve dataset.

The results show that QAOA produces interpretable binary well selection decisions and competitive optimization behavior. In the Volve-based case, QAOA identified a multi-well development configuration, while both the brute-force solver and simulated annealing returned the mathematically optimal single-well solution. This difference highlights the engineering relevance of near-optimal heuristic solutions in practical field development planning. The study demonstrates the feasibility of integrating quantum-inspired optimization into petroleum engineering workflows and provides a foundation for future application to larger real-field datasets.
\end{abstract}

\noindent\textbf{Keywords:} QUBO; QAOA; Oil Reservoir Optimization; Well Placement; Quantum Computing; Hybrid Optimization; Volve Field

\section{Introduction}
Well placement optimization is a central problem in petroleum engineering and field development planning. Selecting the appropriate subset of production wells influences total recovery, sweep efficiency, pressure support strategy, and project economics. However, the search space of possible well combinations grows exponentially with the number of candidate locations, which makes the problem computationally challenging \citep{bangerth2006, yeten2003, onwunalu2010}.

Classical approaches to this problem include exhaustive search, gradient-based procedures, metaheuristics, and simulation-assisted optimization \citep{yeten2003, bangerth2006, onwunalu2010, forouzanfar2013}. Exact solvers provide mathematically optimal solutions for small instances, but their computational cost increases rapidly as the decision space grows. Heuristic techniques such as simulated annealing and related metaheuristics provide more scalable alternatives, though they do not guarantee global optimality \citep{kirkpatrick1983}. In reservoir optimization, such heuristic methods have been widely used because they offer a practical balance between solution quality and computational expense \citep{onwunalu2010, forouzanfar2013, siddiqi2015}.

Recent developments in quantum computing have introduced new approaches to combinatorial optimization. In particular, Ising and Quadratic Unconstrained Binary Optimization (QUBO) formulations provide a convenient bridge between engineering decision models and quantum-inspired optimization frameworks \citep{lucas2014, glover2019}. Among these, the Quantum Approximate Optimization Algorithm (QAOA) has emerged as one of the most promising hybrid quantum-classical methods for discrete optimization \citep{farhi2014, mcclean2016, harrigan2021}. The algorithm operates through a variational loop in which a parameterized quantum circuit is repeatedly evaluated while a classical optimizer updates the control parameters.

This study bridges quantum computing and petroleum engineering by formulating the well placement problem as a QUBO model and solving it using QAOA. The framework is first validated on synthetic benchmark cases and then applied to a realistic field-inspired example derived from the publicly available Volve dataset \citep{equinor2018}. The results are compared with brute-force optimization and simulated annealing in order to evaluate both the mathematical quality and the engineering relevance of the solutions.

The main contributions of this work are:
\begin{itemize}
    \item Formulation of the well placement problem as a QUBO model;
    \item Application of QAOA as a quantum-inspired optimization framework for reservoir decision-making;
    \item Benchmarking against brute-force search, simulated annealing, and random search;
    \item Extension to a realistic field-inspired case using Volve-based data.
\end{itemize}

\section{Methodology}

\subsection{Binary Decision Structure}
Let each candidate well location be represented by a binary decision variable
\[
x_i \in \{0,1\},
\]
where \(x_i = 1\) means that well \(i\) is selected, and \(x_i = 0\) means it is not selected. The optimization seeks a subset of wells that balances production benefit against interference or spacing penalties. Such binary formulations are natural for QUBO modeling and are widely used in combinatorial optimization \citep{lucas2014, glover2019}.

\subsection{QUBO Formulation}
The optimization objective is written in quadratic binary form as
\[
E(\mathbf{x}) = \mathbf{x}^T Q \mathbf{x},
\]
where:
\begin{itemize}
    \item \(Q\) is the QUBO matrix,
    \item diagonal terms represent production reward and cost contribution,
    \item off-diagonal terms represent pairwise well interference.
\end{itemize}
In the simplified structure adopted here,
\[
Q_{ii} = -P_i + C_i,
\]
\[
Q_{ij} = I_{ij}, \quad i \neq j,
\]
where:
\begin{itemize}
    \item \(P_i\) is the normalized production contribution of well \(i\),
    \item \(C_i\) is a cost-related penalty term,
    \item \(I_{ij}\) is the pairwise interference penalty between wells \(i\) and \(j\).
\end{itemize}
This structure captures the engineering trade-off between maximizing productivity and limiting unfavorable inter-well interactions. QUBO formulations are especially useful because they map naturally into Ising-type Hamiltonians suitable for quantum-inspired solvers \citep{lucas2014, glover2019}.

\subsection{QAOA Solution Procedure}
The QUBO matrix is transformed into an Ising Hamiltonian and optimized using QAOA \citep{farhi2014}. The algorithm alternates between a cost unitary based on the problem Hamiltonian and a mixing unitary that explores the feasible binary search space. A classical optimizer updates the variational angles \((\gamma, \beta)\), forming the standard hybrid loop of variational quantum algorithms \citep{mcclean2016, preskill2018}. In the present study, QAOA is implemented in simulation form and used as a quantum-inspired optimizer rather than as a hardware-specific demonstration.

\subsection{Classical Benchmarking}
To evaluate the proposed framework, the following methods are used for comparison:
\begin{itemize}
    \item brute-force exact enumeration,
    \item simulated annealing (SA),
    \item random search,
    \item QAOA.
\end{itemize}
Simulated annealing is selected as a baseline because it is one of the most established metaheuristic optimization methods and has been widely applied in engineering optimization \citep{kirkpatrick1983}. Brute-force search is used only for small cases in order to provide exact reference solutions.

\section{Synthetic Benchmark Results}
Synthetic test cases were evaluated for candidate sets of size \(N = 6,8,10\). Four performance metrics were recorded: objective energy, production indicator, interference penalty, and runtime.

\begin{figure}[H]
    \centering
    \includegraphics[width=0.8\textwidth]{energy_comparison.png}
    \caption{Objective energy comparison across optimization methods for synthetic benchmark cases. Lower values indicate better objective performance.}
    \label{fig:energy}
\end{figure}

\begin{figure}[H]
    \centering
    \includegraphics[width=0.8\textwidth]{production_comparison.png}
    \caption{Production comparison across optimization methods for synthetic benchmark cases. Higher values indicate better production performance.}
    \label{fig:production}
\end{figure}

\begin{figure}[H]
    \centering
    \includegraphics[width=0.8\textwidth]{interference_comparison.png}
    \caption{Interference penalty comparison across optimization methods for synthetic benchmark cases. Lower values indicate lower inter-well interaction penalties.}
    \label{fig:interference}
\end{figure}

\begin{figure}[H]
    \centering
    \includegraphics[width=0.8\textwidth]{runtime_comparison.png}
    \caption{Runtime comparison across optimization methods for synthetic benchmark cases. Lower runtime indicates faster execution.}
    \label{fig:runtime}
\end{figure}

The synthetic benchmark confirms that QAOA is capable of producing competitive objective values relative to classical heuristic methods while maintaining a fully interpretable binary decision structure. This is consistent with the broader view that hybrid quantum-classical methods may become useful in engineering-scale combinatorial optimization \citep{farhi2014, harrigan2021}.

\section{Field-Inspired Case Study Using Volve Data}
To make the framework more realistic, the optimization workflow was extended to a simplified field-inspired case derived from the publicly available Volve dataset \citep{equinor2018}. A subset of five wells was used:
\[
[A1, A2, A3, A4, A5].
\]
The production term was obtained by normalizing average oil rates, while the interference term was derived from pairwise spatial distances using an exponential decay model. This yielded a \(5 \times 5\) QUBO matrix representing the trade-off between productivity and inter-well interference.

\begin{figure}[H]
    \centering
    \includegraphics[width=0.8\textwidth]{production_scatter.png}
    \caption{Production values and selection status for individual wells in the Volve-based case.}
    \label{fig:scatter}
\end{figure}

\subsection{QAOA Results on the Volve-Based Case}
The QAOA run produced the following solution:
\begin{itemize}
    \item Best bitstring: 10110
    \item Binary decision vector: \([1,0,1,1,0]\)
    \item Selected wells: A1, A3, A4
    \item QUBO value of selected solution: 1.136688
\end{itemize}

\subsection{Exact Brute-Force Reference}
The exact brute-force solver returned:
\begin{itemize}
    \item Best bitstring: 00100
    \item Binary decision vector: \([0,0,1,0,0]\)
    \item Selected well: A3
    \item Best QUBO value: -0.900000
\end{itemize}

\subsection{Discussion of the Difference}
The difference between the QAOA and exact results is important from both mathematical and engineering perspectives.

From a strict optimization viewpoint, the brute-force result is globally optimal because it yields the minimum QUBO value. However, it leads to a conservative development strategy in which only one well is selected.

In contrast, the QAOA solution proposes a multi-well development configuration (A1, A3, A4). Although this configuration is not mathematically optimal in QUBO value, it is operationally meaningful because it reflects a broader trade-off between production expansion and tolerable interference.

This behavior is significant in practical field development. In real reservoir management, operators often evaluate near-optimal multi-well strategies rather than adopting the most conservative mathematical minimum, particularly when additional wells may increase cumulative production or improve development flexibility \citep{onwunalu2010, forouzanfar2013, chen2017}.

\subsection{Engineering Interpretation}
The Volve-based case suggests the following:
\begin{itemize}
    \item the framework can translate field-style well data into a QUBO model;
    \item QAOA can generate interpretable well deployment decisions on realistic inputs;
    \item the selected QAOA configuration may represent a development-oriented solution rather than a purely conservative optimum;
    \item the difference between exact and heuristic results is itself useful, as it reveals the engineering trade-off between optimality and operational flexibility.
\end{itemize}

\section{Comparison with Simulated Annealing}
To further evaluate the proposed framework, a simulated annealing algorithm was applied to the same Volve-based QUBO model.

The SA results are summarized as follows:
\begin{itemize}
    \item Best bitstring: 00100
    \item Binary decision vector: \([0,0,1,0,0]\)
    \item Selected well: A3
    \item QUBO value: \(-0.900000\)
\end{itemize}

Interestingly, the simulated annealing solution exactly matches the brute-force optimal result. This confirms that the QUBO formulation is correctly constructed and that classical heuristics are capable of recovering the global optimum for this small-scale problem.

However, the QAOA solution differs significantly:
\begin{itemize}
    \item QAOA selected wells: A1, A3, A4
    \item QAOA QUBO value: 1.136688
\end{itemize}

This discrepancy highlights a fundamental distinction between classical and quantum-inspired optimization approaches. While SA converges toward the global minimum of the mathematical objective, QAOA explores a broader solution landscape and identifies alternative configurations that may be more relevant from an engineering perspective \citep{farhi2014, mcclean2016, harrigan2021}.

\subsection{Engineering Interpretation of SA versus QAOA}
The SA and brute-force solutions correspond to a conservative strategy that selects only the highest-performing well. In contrast, QAOA proposes a multi-well configuration that balances production gain and interference penalties.

From a field development perspective, such configurations are often more desirable, as they:
\begin{itemize}
    \item increase total production capacity,
    \item improve reservoir drainage coverage,
    \item provide operational flexibility.
\end{itemize}
Therefore, the QAOA solution can be interpreted as a near-optimal but more practical development strategy rather than a strict mathematical optimum.

\subsection{Comparative Summary}
\begin{table}[H]
    \centering
    \caption{Comparison of optimization methods on the Volve-based QUBO problem.}
    \begin{tabular}{|l|l|c|c|}
        \hline
        \textbf{Method} & \textbf{Selected Wells} & \textbf{Bitstring} & \textbf{QUBO Value} \\
        \hline
        Brute Force & A3 & 00100 & -0.900000 \\
        Simulated Annealing & A3 & 00100 & -0.900000 \\
        QAOA & A1, A3, A4 & 10110 & 1.136688 \\
        \hline
    \end{tabular}
    \label{tab:comparison}
\end{table}

\section{Conclusion}
This study presented a hybrid quantum-classical framework for oil reservoir well placement optimization using QUBO modeling and QAOA. The framework was first validated on synthetic benchmark cases and then extended to a field-inspired dataset derived from Volve.

The synthetic benchmarks demonstrated competitive optimization behavior for QAOA relative to classical heuristics. The Volve-based case confirmed that the framework can process realistic well-level information and produce actionable binary decisions.

Although the brute-force solver identified the global optimum for the small Volve case, the QAOA solution provided an alternative multi-well strategy that may be more relevant in practical field development planning. The simulated annealing solution matched the exact optimum, which further validated the QUBO construction and provided a useful classical baseline.

These findings support the broader view that quantum-inspired methods may contribute operationally meaningful solutions even when they do not exactly reproduce the mathematical optimum. Future work will focus on:
\begin{itemize}
    \item incorporating larger subsets of real field data,
    \item refining the economic objective function,
    \item integrating additional operational constraints,
    \item extending the framework toward more realistic reservoir development scenarios,
    \item exploring future deployment on emerging quantum hardware platforms.
\end{itemize}

\section*{Code and Data Availability}
The source code and supporting data used in this study are publicly available at: \url{https://github.com/sumerian29/quantum-reservoir-optimizer}

\bibliographystyle{plainnat}
\begin{thebibliography}{99}

\bibitem[Farhi et al.(2014)]{farhi2014}
Farhi, E., Goldstone, J., \& Gutmann, S. (2014). A quantum approximate optimization algorithm. \textit{arXiv preprint arXiv:1411.4028}.

\bibitem[Lucas(2014)]{lucas2014}
Lucas, A. (2014). Ising formulations of many NP problems. \textit{Frontiers in Physics}, 2, 5.

\bibitem[Glover et al.(2019)]{glover2019}
Glover, F., Kochenberger, G., \& Du, Y. (2019). A tutorial on formulating and using QUBO models. \textit{arXiv preprint arXiv:1811.11538}.

\bibitem[McClean et al.(2016)]{mcclean2016}
McClean, J. R., Romero, J., Babbush, R., \& Aspuru-Guzik, A. (2016). The theory of variational hybrid quantum-classical algorithms. \textit{New Journal of Physics}, 18(2), 023023.

\bibitem[Johnson et al.(2011)]{johnson2011}
Johnson, M. W., Amin, M. H. S., Gildert, S., Lanting, T., Hamze, F., Dickson, N., Harris, R., Berkley, A. J., Johansson, J., Bunyk, P., Chapple, E. M., Enderud, C., Hilton, J. P., Karimi, K., Ladizinsky, E., Ladizinsky, N., Oh, T., Perminov, I., Rich, C., Thom, M. C., Tolkacheva, E., Truncik, C. J. S., Uchaikin, S., Wang, J., Wilson, B., \& Rose, G. (2011). Quantum annealing with manufactured spins. \textit{Nature}, 473, 194-198.

\bibitem[Kirkpatrick et al.(1983)]{kirkpatrick1983}
Kirkpatrick, S., Gelatt, C. D., \& Vecchi, M. P. (1983). Optimization by simulated annealing. \textit{Science}, 220(4598), 671-680.

\bibitem[Bangerth et al.(2006)]{bangerth2006}
Bangerth, W., Klie, H., Wheeler, M. F., Stoffa, P. L., \& Sen, M. K. (2006). On optimization algorithms for the reservoir oil well placement problem. \textit{Computational Geosciences}, 10, 303-319.

\bibitem[Yeten et al.(2003)]{yeten2003}
Yeten, B., Durlofsky, L. J., Aziz, K., \& Kazemi, H. (2003). Optimization of nonconventional well type, location, and trajectory. \textit{SPE Journal}, 8(3), 200-210.

\bibitem[Onwunalu \& Durlofsky(2010)]{onwunalu2010}
Onwunalu, J. E., \& Durlofsky, L. J. (2010). Application of a particle swarm optimization algorithm for determining optimum well location and type. \textit{Computational Geosciences}, 14, 183-198.

\bibitem[Forouzanfar et al.(2013)]{forouzanfar2013}
Forouzanfar, F., Li, G., \& Reynolds, A. C. (2013). A two-stage well-placement optimization method based on adjoint-gradient and genetic algorithm. \textit{Journal of Petroleum Science and Engineering}, 112, 254-264.

\bibitem[Siddiqi \& de Souza(2015)]{siddiqi2015}
Siddiqi, A. H., \& de Souza, A. (2015). Reservoir optimization under uncertainty: a review. \textit{SPE Journal}, 20(2), 1-15.

\bibitem[Chen et al.(2017)]{chen2017}
Chen, Y., Oliver, D. S., \& Zhang, D. (2017). Efficient ensemble-based closed-loop production optimization. \textit{SPE Journal}, 22(6), 1737-1756.

\bibitem[Equinor(2018)]{equinor2018}
Equinor (2018). Volve field data release. Available from the Equinor open data portal.

\bibitem[Preskill(2018)]{preskill2018}
Preskill, J. (2018). Quantum computing in the NISQ era and beyond. \textit{Quantum}, 2, 79.

\bibitem[Harrigan et al.(2021)]{harrigan2021}
Harrigan, M. P., Sung, K. J., Neeley, M., Satzinger, K. J., Arute, F., Arya, K., Atalaya, J., Bardin, J. C., Basso, J., Boixo, S., Brooks, P., Brusse, J., Buckley, B. B., Bushnell, N., Chiaro, B., Collins, R., Courtney, W., Demura, S., Dunsworth, A., Eppens, D., Erickson, C., Farhi, E., Fowler, A., Foxen, B., Giustina, M., Graff, R., Habegger, S., Harr, R., Hilton, J., Ho, A., Huang, T., Huggins, W. J., Ioffe, L. B., Jeffrey, E., Jiang, Z., Kelly, J., Klimov, P. V., Knysh, S., Korotkov, A., Kostritsa, F., Landhuis, D., Laptev, P., Lindmark, M., Lucero, E., Martin, O., McClean, J. R., Mi, X., Mohseni, M., Mutus, J., Naaman, O., Neill, C., O'Brien, T. E., Oprencak, A., Petukhov, A., Quintana, C., Rieffel, E. G., Roushan, P., Rubin, N. C., Sank, D., Satzinger, K. J., Schuster, T., Smelyanskiy, V., Syed, M., Vainsencher, A., Villalonga, B., White, T., Yao, Z. J., Yeh, P., \& Martinis, J. M. (2021). Quantum approximate optimization of non-planar graph problems on a planar superconducting processor. \textit{Nature Physics}, 17, 332-336.

\end{thebibliography}

\end{document}